\chapter{Device and execution model}\label{ch:device-execution-model}

All measurements were made on one Apple M5 Pro GPU (G17, the 20-core H17s) under the software listed in
\cref{sec:measured-part-software}. Later chapters assume the execution model set out here. Mechanisms are in \crefrange{ch:register-file}{ch:kernel-metadata-image} and measured costs in \cref{ch:performance}.

These terms are used throughout the document:

\begin{itemize}
\item \emph{Lane}: one thread's slot in a simdgroup. \emph{Simdgroup}: 32 lanes that issue one instruction stream together under an
  execution mask (\cref{ch:control-flow-execution}). \emph{Threadgroup}: one or more simdgroups that share threadgroup memory and can meet at a
  threadgroup barrier (\cref{sec:barriers-fences}). \emph{Grid}: all threadgroups of one dispatch.
\item \emph{Core}: one of the 20 GPU cores. Per-core figures divide a GPU-wide count by 20.
\item \emph{Resident}: a threadgroup (or simdgroup) is resident while it holds a core's resources, from launch until it ends.
  \emph{Residency} or \emph{occupancy}: how many are resident at once.
\item \emph{Registers used}: the highest 32-bit register a program names, plus one. \emph{Declared register count}: the value in
  per-kernel metadata slot 0 (\cref{sec:gpu-metadata-slots}).
\item \emph{Warm} and \emph{cold} refer to the GPU's performance state (\cref{sec:clocks-performance-states}).
\item "MM x" cites the project's machine model, \code{docs/g17-tensorops-machine-model.md}.
\end{itemize}

\section{Platform versions}\label{sec:measured-part-software}

The measurements were taken on one machine and one toolchain, on the platform in \cref{tab:measured-platform} (MM~\mm{22}, MM~\mm{25.66},
MM~\mm{25.91}). The below-Metal evidence (\cref{ch:submission-below-metal}) records the OS build, 25G83; the Metal-path
evidence files do not, so for those measurements the build is the stated platform rather than a per-file record.
On 2026-10-02 the machine was updated to macOS 27.0.1 (build 26A434). Nothing in this reference was re-measured on that
build; \code{tools/g17platform.py --check} reports the change, and the table gives the platform the figures were taken on.
The update renumbered Apple's decoder; the tools translate back to the numbering used here (\cref{sec:decoder-renumbering}).

\begin{xltabular}{\linewidth}{@{}>{\hsize=0.65\hsize}X>{\hsize=1.40\hsize}XX@{}}
\caption{The measured platform: machine, OS, driver and toolchain versions, with how each was read.}\label{tab:measured-platform}\\
\toprule
Component & Value & How read \\
\midrule
\endfirsthead
\tablecontinued{3}
\toprule
Component & Value & How read \\
\midrule
\endhead
Machine & \code{Mac17,8}, Apple M5 Pro & \code{sysctl hw.model} \\*
GPU & "Apple M5 Pro", 20 cores, Metal 4; Apple's internal
names H17s, \code{applegpu\_g17s} & \code{system\_profiler}; the core count is the
campaign's own record (below) \\*
macOS & 26.6.2, build 25G83; Darwin 25.6.0 & \code{sw\_vers}, \code{uname} \\*
GPU user driver & \code{AGXMetalG17X.bundle} 353.14 & MM~\mm{25.141.20} \\*
Metal compiler & metal 32023.921 (metalfe-32023.921.6), target air64-apple-darwin25.6.0,
against an OS-side GPU compiler at 32023.886 & \code{xcrun metal --version}; MM~\mm{22} \\*
Xcode / SDK & Xcode 27.0 (27A266a), macOS SDK 27.0 & \code{xcodebuild -version} \\*
Host compiler & Apple clang 21.0.0 (clang-2100.3.34.2) & \code{clang --version} \\*
Apple's AGX3 decoder & inside
\code{GPUCompiler.framework/Versions/32023/Libraries/libLLVM.dylib} & \cref{sec:decoder} \\*
\bottomrule
\end{xltabular}

\code{sysctl -n hw.ncpu} returns 18, the CPU core count, which the project has mistaken for the GPU core count before. The GPU core count is not
readable through sysctl and is carried as the campaign's own record; a platform pin that confused the two would put a
wrong divisor under every per-core figure (MM~\mm{25.66}). \code{tools/g17platform.py --check} compares the running machine with
this record and names each field that differs, with what it bears on: the compiler version on lowering, scheduling
and every "Apple's compiler emits" statement; the OS build on driver behaviour and loader variants; the part on
everything (MM~\mm{25.66}).

Results on this part do not establish results on any other part,
OS build or compiler: another G17 product (a different core count, a different populated register file, other
loader variants) is a different population, and no row of the project's ledger claims otherwise (MM~\mm{0.1}, MM~\mm{25.66}
row T9, MM~\mm{25.119}). MM~\mm{25.119} groups results by how likely
they are to change on another part, OS build or compiler:

\begin{itemize}
\item \emph{Properties of the toolchain}, which follow Apple's compiler rather than the silicon and are the likeliest to move:
  which source forms reach the tensor unit, metadata layout classes, the constant-program laws, which opcodes Apple's
  compiler selects, spill counts (MM~\mm{25.119}).
\item \emph{Properties of the core count and clock}: residency counts, speeds, power (MM~\mm{25.119}).
\item \emph{Encodings and arithmetic}, the least likely to move within this generation: field maps, the register namespace,
  the tensor unit's reduction arithmetic, which "has held across every G17 measurement in this repository" (MM~\mm{25.119}).
  Several conventions are conserved from Apple's earlier G13 generation and its public reconstruction in Mesa: the
  length-bit cluster in the encoding (\cref{sec:framing}) and the command-stream layout below Metal, which decodes field for
  field against Mesa's G13 description (\cref{ch:submission-below-metal}). Mesa's model was used only to suggest hypotheses and is not evidence for G17.
\item Apple's decoder admits the same ten tensor MMA opcodes when told the CPU is g17s, g17g or g18, and decodes the same
  bytes as unrelated instructions for g17, g16 and g15 targets (MM~\mm{2}). This is a property of Apple's decoder tables
  and is no evidence about other silicon.
\end{itemize}

The portable parts of the project's own work (byte-identity pins, static checks, the metadata determinant over the
committed corpus) rerun anywhere on the CPU; the platform facts each carry a named re-measurement instrument
(MM~\mm{25.119}).

\section{Threads, simdgroups, threadgroups and grids}\label{sec:threads-simdgroups-threadgroups}

\begin{itemize}
\item SIMD width is 32 (measured; \code{threadExecutionWidth} is the driver constant 32 for every pipeline, MM~\mm{25.141.20},
  MM~\mm{25.147}). A tensor fragment is a per-lane register tuple across those 32 lanes (\cref{ch:tensor-unit}).
\item A threadgroup holds at most 1,024 threads. \code{maxTotalThreadsPerThreadgroup} is a field of the compiler's reply that
  the driver copies; 0 or absent becomes 1,024 (driver, read statically, MM~\mm{25.141.20}). It read 1,024 for every pipeline
  tested, from 4 to 512 live scalars per thread and 64 to 32,768 bytes of threadgroup memory (measured, MM~\mm{25.47}), so it
  does not vary with register use.
\item Registers are per lane and never cross simdgroups. Every hand-off between simdgroups goes through memory and a
  barrier (\cref{sec:barriers-fences}); cross-lane movement inside a simdgroup is \cref{ch:cross-lane-operations}.
\item Threadgroup memory is limited to 32,768 bytes per threadgroup, a hard limit enforced by Apple's GPU compiler on this toolchain (48, 60
  and 64 KiB refused; 96 KiB drops the compiler connection; MM~\mm{7}, MM~\mm{0.1}). Metal's pipeline creation does not check the
  declared size: a container declaring 1 MiB builds and reports 1,048,576 back (MM~\mm{25.147}; \cref{sec:object-library-archive}).
\item Position registers (lane index, simdgroup index, threadgroup and grid positions) are special registers read by
  instruction; \cref{ch:special-registers} is the table.
\item Grids are dispatched in threadgroup order, and the driver prepares fast-division constants to split a linear
  threadgroup index into three dimensions (MM~\mm{25.141.20}). The ordering consequence is \cref{sec:forward-progress-in-order}.
\end{itemize}

\section{Core resource model}\label{sec:core-logical-resource}

The model in \cref{tab:core-resources} is a \emph{logical} account of measured concurrency: which streams run side by side without slowing each
other, and which contend. It is not a floorplan and does not locate or count physical units (MM~\mm{0.2}).

\begin{xltabular}{\linewidth}{@{}>{\hsize=0.65\hsize}X>{\hsize=1.59\hsize}X>{\hsize=0.75\hsize}X@{}}
\caption{Logical resources of one core and their measured concurrency behaviour.}\label{tab:core-resources}\\
\toprule
Resource & Measured behaviour & Source \\
\midrule
\endfirsthead
\tablecontinued{3}
\toprule
Resource & Measured behaviour & Source \\
\midrule
\endhead
Four SIMD issue slots per core & One, two or four simdgroups each running a dependent tensor MMA chain
take 32.8, 32.4 and 32.8~µs; eight take 57.5, sixteen 112.6, thirty-two
224.3: flat to four, linear beyond & measured, MM~\mm{8} \\*
Tensor unit, one slot per SIMD & 16 × 16 × 16 per issue, 1,024 flop per cycle per core at saturation;
about 5.5 simdgroups per core saturate it & measured, MM~\mm{8}; \cref{ch:tensor-unit} \\*
fp32 / integer ALU & about 4.1 × $10^{12}$ fp32 adds per second machine-wide, consistent with
128 fp32 lanes per core at about 1.6~GHz; no dual issue between
fp32 and int32 (ffma + iadd mixed 3.61~ns per op against 3.64 predicted
for one shared port and 2.09 for dual issue) & measured, MM~\mm{0.15} \\*
Legacy simdgroup-matrix engine & overlaps the tensor unit completely at 16 and 32 simdgroups per core;
its fp32-accumulate forms contend with fp32 FMA (overlap 0.37 or less),
its half-accumulate form mostly does not (0.66–0.95) & measured, MM~\mm{9}; \cref{sec:legacy-simdgroup-matrix-engine} \\*
Per-simdgroup scoreboard & eight slots ordering late values (loads, special-register reads,
atomics) against their consumers & measured, MM~\mm{6}; \cref{sec:scoreboard} \\*
Load/store path & per-lane addresses; no hidden tile-layout unit & measured, MM~\mm{23.2}; \cref{ch:memory} \\*
\bottomrule
\end{xltabular}

"Contends" and "overlaps" in this table are exchange rates measured at stated
occupancies; what physically is shared was not identified ("about 8 fma of budget per legacy issue" is a rate, not a
mechanism; MM~\mm{9}). Partial overlap at one or two simdgroups per core is the in-order issue of a single
simdgroup, not a shared unit (MM~\mm{9}). Performance figures are in \cref{ch:performance}.

\section{Residency and occupancy rules}\label{sec:residency-occupancy-rules}

These are the rules a compiler or runtime may rely on; the measured residency counts, the per-core staircase and what
registers and spills cost are in \cref{sec:registers-spills-occupancy,sec:residency-occupancy-measurements}.

\begin{enumerate}
\def\labelenumi{\arabic{enumi}.}
\item Residency follows the registers a program uses, not the count it declares (measured, MM~\mm{25.115}). The same bytes
  declaring 19 and declaring 125 reach the same peak every time (1,814–1,896 resident at 2,048 launched simdgroups); only
  using more registers lowers the peak. The declared count (metadata slot 0) is inert for residency between 19 and
  125, and inert for execution: code using R0–R123 runs correctly declaring 124, 127, 128, 160 or 224 (MM~\mm{10}).
  Below Metal the count travels to the hardware in a launch-state packet (\cref{sec:gpu-metadata-slots}, \cref{ch:submission-below-metal}); whether that
  packet's value affects residency was not isolated.
\item The driver caps threadgroups per core at \code{min(3072 / threads\_per\_threadgroup, 96)} (static read of the driver,
  MM~\mm{25.141.20}): 3,072 threads per core, 96 simdgroups of 32. Nothing in that path reads the register count; any
  register-dependent limit sits below the driver or in the compiler's reply. The measured peak at 19 registers,
  95.0 simdgroups per core (MM~\mm{25.121}), sits at this cap.
\item Register pressure lowers residency, by a kernel-dependent amount. Two probes disagree by about a factor of two
  at high register counts (about 64 against about 32 simdgroups per core near 125 registers; the counts and their
  kernels are \cref{sec:residency-occupancy-measurements}), and which kernel property the difference follows is open (MM~\mm{25.121}). A scheduler must use
  the count measured for its own kernel family.
\item The per-core threadgroup-memory budget is 57,344 bytes (56 KiB) (measured, MM~\mm{25.49}): two threadgroups of
  28,672 bytes fit on a core, one of 28,688 does, and the 16-byte step costs 25 percent. With the 32 KiB per-threadgroup
  limit, at most one threadgroup above 28 KiB is resident per core.
\item For 1,024-thread threadgroups, no residency count may be assumed. The forward-progress probe ran 200 threadgroups of
  1,024 threads, one producer last, with no consumer starving (17–21~µs); had fewer than 200 been co-resident, the
  earliest consumers would have spun to the 200,000-poll cap before the last-dispatched producer ran (MM~\mm{25.141.1}).
  That implies at least 10 per core. The driver formula gives 3 per core (60 GPU-wide) (MM~\mm{25.141.20}). The two numbers
  have not been reconciled, and the 800-threadgroup row (480 of 799 at the cap) does not give a simultaneous count,
  because waves can reach the cap one after another. Do not size a spin-waiting grid of 1,024-thread groups from
  either number (\cref{sec:forward-progress-in-order}).
\end{enumerate}

\textbf{Not known.} The physical register capacity per core and its allocation granularity. Residency is not one fixed
pool divided by per-simdgroup demand (a fixed-pool model under-predicts the measured count about fourfold at 126
registers; \cref{sec:residency-occupancy-measurements}, MM~\mm{24.2} row H4, MM~\mm{25.49}), and whether that means a larger file than any published part or
registers that are not all physically resident is not separable with these instruments (MM~\mm{25.20}). Metal's public API
has no occupancy counter on this part (MM~\mm{25.48}; \cref{sec:dispatch-encoder-command-buffer}).

\section{Forward progress}\label{sec:forward-progress-in-order}

The probe is one dispatch (\code{tools/g17gpuprobes/fp.m}): consumers spin, capped at 200,000 polls, until an atomic count
reaches the number of producers placed at the end of the grid (MM~\mm{25.141.1}). \Cref{tab:forward-progress-probe} gives its results.

\begin{xltabular}{\linewidth}{@{}XlX@{}}
\caption{Forward-progress probe results by grid shape, with consumers capped at 200,000 polls (MM~\mm{25.141.1}).}\label{tab:forward-progress-probe}\\
\toprule
Grid (threadgroups × threads) & Producers & Consumers at the cap \\
\midrule
\endfirsthead
\tablecontinued{3}
\toprule
Grid (threadgroups × threads) & Producers & Consumers at the cap \\
\midrule
\endhead
64 to 1,200 × 32 & 1, the last & 0 (none waited more than 94 polls) \\*
2,000 × 32 & 1 & 1,688 of 1,999 \\*
20,000 × 32 & 1 & 18,350 of 19,999 \\*
20 to 200 × 1,024 & 1 & 0 \\*
800 × 1,024 & 1 & 480 of 799 \\*
\bottomrule
\end{xltabular}

These are observations of one probe family in the tested configurations, not a scheduling guarantee. Below, each is
used only to restrict what code may assume.

\begin{itemize}
\item Among resident threadgroups, no starvation was observed. Up to 1,200 single-simdgroup threadgroups and up to 200
  of 1,024 threads, every consumer saw its producer within the cap (at most 94 polls in the 32-thread rows).
\item Past residency, waiters starved. With 2,000 threadgroups the first resident wave (about 1,690, about 84 per core at
  that register count) spun to its cap before the producer, dispatched last, ran, and the threadgroups that
  succeeded were the high-numbered ones. That fits in-order dispatch of threadgroups with no preemption of a
  spinning one in these runs (MM~\mm{25.141.1}); no run tested whether preemption can occur under other
  conditions, so code must not rely on preemption to make progress.
\item In the persistent-kernel runs (16 phases, totals exact), grid barriers among resident threadgroups all completed
  without reaching the cap, and a barrier cost more than a dispatch boundary at every measured shape (MM~\mm{25.141.1};
  costs in \cref{ch:performance}). A grid that waits on itself has to be sized to resident capacity at its own
  register and threadgroup-memory use, measured for that kernel (\cref{sec:residency-occupancy-rules}).
\end{itemize}

The following may not be assumed:

\begin{itemize}
\item that a threadgroup waiting on another threadgroup will see it run, unless both are resident; a spin-wait on a
  producer that is not resident may wait indefinitely. In the probes, polls were capped, and waiters beyond residency
  spun to the cap without seeing their producer run (MM~\mm{25.141.1}; MM~\mm{17} rule 40, MM~\mm{0.15});
\item that "in order" means anything finer than dispatch order of threadgroups: completion order, and ordering of memory
  effects between threadgroups, are separate questions answered in \cref{sec:ordering-across-lanes} (visibility needs a device fence on both
  sides; ordering comes from an atomic);
\item that a timeout protects the machine. Killing the host process does not cancel a kernel: an unbounded loop kept the GPU
  at 100 percent with nothing dispatching until a reboot (MM~\mm{25.116}), and an unbounded data-driven trip count once
  rebooted the machine (MM~\mm{25.90}). Loops in this project's compiler are capped so that termination is provable
  (\cref{sec:loops-loop-variable-rule}, \cref{ch:control-flow-execution}).
\end{itemize}

Within those configurations, the measurements support a persistent kernel that sizes its grid to resident capacity at
its own register and threadgroup-memory use and pulls work from an atomic queue; any other hand-off between
threadgroups goes through a dispatch boundary (MM~\mm{17} rule 40).

\section{Clock states and timing}\label{sec:clocks-performance-states}

Timing on this GPU follows the rules below; the measurements are in \cref{sec:clocks-power-gpu}.

\begin{itemize}
\item The top clock is about 1.62~GHz: 1,620~MHz read from \code{powermetrics} at 100 percent residency in the top bin;
  1.63~GHz is implied by the saturated tensor issue rate of 32 cycles per MMA (measured, MM~\mm{8}, MM~\mm{20}).
\item A run shorter than about 1.5~ms runs at whatever frequency the GPU is
  idling at, and a longer one speeds up partway (MM~\mm{25.121}); a cold GPU ran the same kernel about 4 to 4.8 times slower
  than after sustained load (MM~\mm{22}, MM~\mm{25.115}). Timings need a saturating warm-up in the same process (the project uses
  4,096 threadgroups of 256) and a reference dispatch to detect a clock change.
\item The governor reacts to command-buffer shape. Metal System Trace records the GPU performance state (Minimum,
  Medium, Maximum) per interval; a workload issuing seven short command buffers per step spent 27–78 percent of its time
  in Medium, while one long command buffer per step stayed at Maximum (two recordings, direction reproduced, size not;
  MM~\mm{0.15}).
\item In-run variation was 0.4–1.3 percent, but the same configuration in two
  processes differed by 9–13 percent; compare arms interleaved in one process (MM~\mm{16}, MM~\mm{25.126}; further figures at the start of \cref{ch:performance}).
\item Several results published earlier in the project were artefacts of the idle clock (a
  residency "shortfall", a "bimodal" issue rate); each is corrected where it is used (MM~\mm{25.115}, MM~\mm{25.122}).
\end{itemize}

\textbf{Evidence.} MM~\mm{0.1}, MM~\mm{0.2}, MM~\mm{0.15}, MM~\mm{2}, MM~\mm{7}, MM~\mm{8}, MM~\mm{9}, MM~\mm{10}, MM~\mm{16}, MM~\mm{17} (rule 40), \mm{20}, \mm{22}, \mm{24.2} (row H4),
\mm{25.20}, \mm{25.47}, \mm{25.48}, \mm{25.49}, \mm{25.66}, \mm{25.90}, \mm{25.91}, \mm{25.115}, \mm{25.116}, \mm{25.119}, \mm{25.121},
\mm{25.122}, \mm{25.126}, \mm{25.141.1}, \mm{25.141.20}, \mm{25.147}.
